\scanpage{11}
De même, on aura (prop. 2, cor)
\[
\varphi_L(xy)=(\varphi_L(x)*\varphi_L(y))\deleted{\unsure{\mathrm{cl}}}\cdot(-1)^n\omega_A
\]
d'où en appliquant $\mathrm{ch}$ et $({*})$
\[
\psi_D(\mathrm{ch}(xy))=\deleted{\unsure{\psi_D(}}\psi_D(\mathrm{ch}\,x)\,\psi_D(\mathrm{ch}\,y))T(-1)^n\deleted{\exp}\mathrm{ch}\,\omega_A
\]
donc comme $\mathrm{ch}(xy)=\mathrm{ch}\,x\ \mathrm{ch}\,y$
\[
\boxed{\psi_D(ab)=\psi_D(a)\,\psi_D(b)\,(-1)^nT\,\mathrm{ch}\,\omega_A}
\quad\left\{\begin{aligned}
T&=\operatorname{Todd}(t_A)\\
\omega_A&=\bigwedge^n t_A^\vee
\end{aligned}\right.
\]
% slanted
\margin{les formules sont valables dès que $a,b\in\mathrm{Gr}^{\cdot}(A)_{\mathbb Q}$.}

Ce sont les formules \deleted{\unsure{affirmées|différentielles}} \inserted{\unsure{démontrées}} (\unsure{resp|resp.} affirmées) par Kleiman dans ``10 Exposés\ldots'', sur un corps alg clos et en travaillant avec une théorie \unsure{cohomologique|chomologique} plutôt qu'avec $\mathrm{Gr}_{\mathbb Q}$, -- à cela près que Kleiman\deleted{\unsure{n}} n'a pas tenu compte du facteur $(-1)^n$ !

4) Il résulte du Th. de dualité pour $A\times_S B\longrightarrow B$ que l'on a un isom fonctoriel
\[
\Phi_L(\check F\otimes\omega_A[n])\simeq\Phi_L(F)^\vee
\]
i.e.
\[
\boxed{\Phi_L(F)^\vee\simeq\Phi_L(\check F\mathbin{\cdot}\omega_A[n]}
\]
Du point de \unsure{vue|vu} virtuel, ceci s'écrit
\[
\varphi_L(x)^\vee=(-1)^n\omega_A\,\varphi_L(\check x)
\]

5) Bien sûr, si $u$ est un autom de $A$ et $u'$ l'autom contragrédient de $B$, on a par transport de structure un isom canonique
\[
\Phi_L(u.F)\simeq u'.\Phi_L(F)
\]
En particulier, $\Phi_L$ commute à la symétrie $s_A$. Kif-kif pour $\varphi_L$.
% archive sheet number: 10
